\section{Failure Mode Analysis}
\paragraph{Instruction: a constraint obscured by overall impressions.}
For an academic talk restricted to 5--10 slides, Manus produces 12, with HCR .50. SlideAudit, PPTBench-style, and LecEval nevertheless score the deck .993, .851, and .873. These implemented scores do not separately test whether the count falls inside the requested range. SlideGate isolates the count constraint and checks it programmatically.

\paragraph{Content: using the wrong source.}
A training task requests materials from lesson 2, but Doubao builds its deck from lesson 1. The deck appears complete and normally formatted, while the requested source is unused. Scoring slide coherence without consulting the source cannot directly reveal this mismatch. Source-grounded coverage identifies the missing items; Appendix~\ref{app:evidence} shows page-level evidence.

\paragraph{Vision: appealing appearance does not establish complete content.}
Kimi has the highest visual mean (.945), while its content interval overlaps those of several other systems. Manus has the highest content mean but a lower visual score (.886). Visual and content correlations with the appearance-focused human ranking are .855 and .029, respectively. Content coverage and visual defects therefore require separate evidence.

\begin{table}[t]
\caption{Representative failure cases. These examples illustrate diagnosis rather than estimate population failure frequencies.}
\label{tab:cases}\centering\small
\begin{tabularx}{\linewidth}{@{}YY Y@{}}
\toprule Task and system & Observed failure & Comparison scores or missing check\\\midrule
Academic talk, 5--10 slides; Manus & 12 slides; HCR .50 & SlideAudit .993; PPTBench-style .851; LecEval .873\\
Edit to 18--22 slides with bracketed placeholders; Manus & 10 slides, no placeholders; HCR .77 & SlideAudit .993\\
Training deck grounded in lesson 2; Doubao & Uses lesson 1; required lesson-2 coverage is low & Source materials are not inspected by the compared content scoring setup\\
Academic talk, 20 frozen content items; AnyGen & Coverage 0; visual score .886, no page defects & No item-level coverage check in the compared scores\\\bottomrule
\end{tabularx}
\end{table}

\section{Discussion}
Fast overall preference captures salient visual attributes but provides little evidence about source-grounded content. In a separate historical aggregation analysis, content-only scores correlate with human preference at .029, equal-weight three-dimensional scores at .429, scores omitting content at .600, and visual-only scores at .829. This historical visual reference differs from the .855 validation result; the two are not interchangeable. Fixed aggregation can dilute the appearance signal in appearance-focused human evaluation while obscuring which dimension requires repair. The same-backbone diagnostic experiment likewise reduces instruction-failure decks from 15 to 1 without a significant visual change.

SlideGate consequently reports dimensions and failure evidence by default. Teaching may emphasize visual clarity, while data analysis may emphasize accurate content. Applications requiring a scalar score should state their weights and also report the individual dimensions.

\section{Limitations}
Visual preference validation concerns system rankings, not a general guarantee of page-level detection accuracy. Programmatic instruction references cover slide count and aspect ratio; other constraints rely on model judgments. Human content annotation concentrates on checklist items, and source faithfulness lacks comparable-scale human gold labels. Checklist consistency has been measured, but completeness and faithfulness have not been systematically validated against an independent human specification. The task collection is small, and complete-case screening may introduce selection bias. Failure labels are not independent human deliverability judgments. Within-dimension weight sensitivity is not established by the cross-dimension analysis. We plan to release task schemas, frozen specifications, and evaluation scripts, with attachments subject to access and redistribution permissions.

\section{Conclusion}
SlideGate combines a presentation-generation benchmark with task-specific verification and explicit failure diagnosis. Specifications are constructed from requests and parseable sources and frozen across systems solving the same task. Among 206 commercial-system outputs with complete comparison scores and determinate failure-union labels, 124 contain at least one detected blocking failure. Existing comparison scores exhibit high false-positive rates when tuned for high failure recall, while explicit checks localize instruction violations, missing content, and visual defects to items and pages. The diagnostic generation comparison also reveals improvements that differ by dimension. We recommend reporting dimensional results and failure evidence, and making threshold-selection procedures and detection--false-alarm tradeoffs explicit when using quality scores for screening.
